\markboth{REFERENCIAS}{REFERENCIAS}  
\addcontentsline{toc}{section}{Referencias}
\renewcommand{\refname}{Referencias}  %<--cambiar ese renglon e os de arriba segun o idioma empregado. Pode chamarse tamen bibiografia no caso de que non haxa notas coma o do bibitem{nota} do noso exemplo

\begin{thebibliography}{99}
%\footnotesize %descomentar para cambiar o tamano de letra desta seccion (footnotesize e valido)
%Exemplo
\bibitem{patricia} J. Abraham et al., ``The fluorescence detector of the Pierre Auger Observatory'', {\it Nuclear Instruments and Methods in Physics Research A}, vol. 620, pp. 227-251, 2010.
\bibitem{nota} \'As veces as entradas bibliogr\'aficas poden conter textos explicativos, como se fosen notas ao p\'e. Tam\'en poden ir varias referencias bibliogr\'aficas nunha soa entrada. Non \'e obrigatorio ning\'un formato espec\'{\i}fico de referencias (orde dos volume-p\'axinas-ano, por exemplo) pos hai moitos en uso internacionalmente. A secci\'on de bibliograf\'{\i}a computa na contaxe de p\'axinas de cara ao l\'{\i}mite m\'aximo da extensi\'on da memoria. No c\'odigo do ficheiro {\tt bibliografia.tex} pode usarse a opci\'on de cambiar tama\~nos, e incluir outras. Lembra compilar varias veces o documento final para que as referencias est\'en correctamente actualizadas.
\end{thebibliography}
